\documentclass[10pt,journal,compsoc]{IEEEtran}

\usepackage{amsmath,amssymb,amsfonts}
\usepackage{algorithmic}
\usepackage{graphicx}
\usepackage{textcomp}
\usepackage{xcolor}
\usepackage{booktabs}
\usepackage{listings}
\usepackage{microtype}
\usepackage{hyperref}
\usepackage{cite}

\lstset{
  basicstyle=\ttfamily\footnotesize,
  keywordstyle=\color{blue}\bfseries,
  commentstyle=\color{gray}\itshape,
  stringstyle=\color{teal},
  breaklines=true,
  frame=single,
  tabsize=2,
  showstringspaces=false
}

\begin{document}

\title{Project Chimera: Empirical Boundaries of Undecidability and Accidental Turing-Completeness in Modern Type Systems}

\author{Lead Compiler Architect \& Formal Type Theorist\\
\IEEEmembership{Formal Systems \& Compiler Architecture Laboratory, Project Chimera}
}

\markboth{IEEE Transactions on Programming Languages and Systems,~Vol.~48, No.~3,~September~2026}%
{Chimera: Empirical Boundaries of Undecidability in Type Systems}

\IEEEtitleabstractindextext{%
\begin{abstract}
Modern production programming language type systems are rarely designed with the overt objective of general-purpose computation. Yet, through the confluence of bounded existential quantification, distributive conditional type narrowing, recursive type aliases, and nominal trait Horn-clause resolution, modern compilers have accidentally achieved Turing-completeness. Because the Halting Problem is strictly undecidable, static compilers cannot mathematically guarantee termination for arbitrary well-formed type expressions without compromising completeness or enforcing operational circuit breakers.

This paper presents \textbf{Project Chimera}, a comprehensive empirical and theoretical study that systematically maps, measures, and stress-tests the boundary where static type checking transitions from polynomial-time verification into super-polynomial explosion, non-termination, and undecidability. We construct a multi-language comparative testbed spanning \textbf{TypeScript 7.0.2}, \textbf{Rust 1.97.0}, and \textbf{Apple Clang C++20}, hosting pure compile-time universal computational substrates including Cook's Turing-complete Rule 110 cellular automaton, Post 2-tag systems, Peano-Ackermann arithmetic, an AST-interpreting Brainfuck virtual machine, and a Kleene Second Recursion Theorem self-reproducing quine.

Our empirical findings reveal: (1) TypeScript enforces a rigid tri-fuse circuit breaker hierarchy: a 48-frame call-stack limit, a 999-iteration tail-call fuel counter, and an absolute global ceiling of $5.03 \times 10^6$ instantiations; (2) By designing a trampolined chunking operator, we bypass the 999-step ceiling to execute $S = 131,072$ steps in 214\,ms; (3) We construct a pathological 25-line Rust trait projection that completely evades Chalk cycle detectors, inducing a 30.4-second clean compiler freeze at depth 23 ($8.38 \times 10^6$ nodes); and (4) In the comparative triad benchmark, Clang C++20 template metaprogramming dominates compile-time execution speed, completing 1,000 steps in 30\,ms ($8.8\times$ faster than Rust, $20.6\times$ faster than TypeScript) and scaling cleanly to 10,000 steps in 470\,ms under $-ftemplate-depth=30000$. We conclude with formal architectural recommendations for future multi-dimensional thermodynamic compiler fuel metering.
\end{abstract}

\begin{IEEEkeywords}
Type Systems, Accidental Turing-Completeness, Undecidability, Compiler Verification, Circuit Breakers, Cellular Automata, Kleene Recursion Theorem.
\end{IEEEkeywords}}

\maketitle
\IEEEdisplaynontitleabstractindextext
\IEEEpeerreviewmaketitle

\section{Introduction}
\IEEEPARstart{S}{tatic} type systems were historically conceived as proof systems to ensure program safety, prevent runtime memory corruption, and guide code generation. Under the Curry-Howard-Lambek correspondence, types correspond to propositions in intuitionistic logic, and program terms correspond to proof normalization trees. In strongly normalizing calculi—such as Girard's System $F$ or Martin-L\"of Type Theory—every well-typed expression terminates unconditionally.

However, over the past two decades, industrial programming languages have prioritized developer expressiveness, complex domain-specific modeling, and type-level reflection. Languages such as TypeScript, Rust, and C++ have enriched their type checkers with advanced constructs:
\begin{itemize}
    \item \textbf{TypeScript:} Distributive conditional types ($T \text{ extends } U ? A : B$), recursive type aliases, variadic tuple spread manipulation, and template literal type decomposition.
    \item \textbf{Rust:} Nominal trait systems with associated types and where-clause unification, mathematically equivalent to first-order Horn-clause logic programming (Prolog) evaluated via SLD-resolution.
    \item \textbf{C++:} Template metaprogramming with non-type template parameters, partial class template specialization, and SFINAE/Concepts constraints.
\end{itemize}

These features inevitably cross the boundary into \textbf{Accidental Turing-Completeness}. By Turing's undecidability theorem (1936) and Rice's theorem (1953), no algorithm can determine whether an arbitrary type computation will terminate:
\begin{equation}
\forall \mathcal{M}, \quad \text{TypeCheck}(\Gamma, \tau) \iff \text{Halts}(\mathcal{M}_\tau)
\end{equation}

Because compilers cannot resolve the Halting Problem, production compilers must deploy heuristic \emph{circuit breakers}—depth limits, fuel counters, and cycle detectors. These mechanisms do not solve undecidability; they merely truncate the execution search tree.

\textbf{Project Chimera} provides the first unified empirical and theoretical characterization of this boundary across three premier modern languages.

\section{Theoretical Foundations and Substrates}

\subsection{Matthew Cook's Rule 110 Universality Substrate}
To study compiler execution without the arbitrary artifacts of ad-hoc benchmarks, we require a minimal, mathematically proven universal computational model. We select Matthew Cook's (2004) landmark theorem: the \textbf{Rule 110 Elementary Cellular Automaton} is universal via the simulation of cyclic tag systems.

Rule 110 maps a 1D tape $c \in \{0, 1\}^W$ across discrete time steps $t \in \mathbb{N}$ via the local neighborhood transition function $f_{110}: \{0, 1\}^3 \to \{0, 1\}$ defined by the byte $01101110_2$:
\begin{align}
f_{110}(1, 1, 1) &= 0, \quad f_{110}(1, 1, 0) = 1, \quad f_{110}(1, 0, 1) = 1, \nonumber \\
f_{110}(1, 0, 0) &= 0, \quad f_{110}(0, 1, 1) = 1, \quad f_{110}(0, 1, 0) = 1, \nonumber \\
f_{110}(0, 0, 1) &= 1, \quad f_{110}(0, 0, 0) = 0
\end{align}

Because $f_{110}(0, 0, 0) = 0$, an infinite sequence of zeros is \emph{quiescent} (time-invariant), ensuring that finite boundary simulations compute without edge corruption.

\subsection{Kleene's Second Recursion Theorem and Type-Level Quines}
Kleene's Second Recursion Theorem (1938) establishes that for every computable function $F$, there exists an index $e$ such that $\phi_e \simeq \phi_{F(e)}$. In program terms, there exists a self-reproducing program $Q$ (a quine) satisfying:
\begin{equation}
\text{eval}(Q) \equiv \text{repr}(Q)
\end{equation}

In type systems, trivial recursion $\text{type } Q = Q$ is detected at bind-time and rejected with circularity errors (e.g., TS2456). A genuine quine must synthesize its abstract syntax tree (AST) dynamically via the diagonal operator:
\begin{equation}
\delta(x) = \text{App}(x, \text{Quote}(x)), \quad Q = \text{App}(\text{Diag}, \text{Quote}(\text{Diag}))
\end{equation}

When evaluated by the type reducer $\text{Resolve}\langle \cdot \rangle$:
\begin{align}
\text{Resolve}\langle Q \rangle &= \text{Resolve}\langle \text{App}(\text{Diag}, \text{Quote}(\text{Diag})) \rangle \nonumber \\
&= \text{App}(\text{Diag}, \text{Quote}(\text{Diag})) = Q
\end{align}

\subsection{Compile-Time Brainfuck Virtual Machine}
To further demonstrate unconstrained Turing universality, we implement a pure compile-time Brainfuck interpreter featuring a bi-infinite zipper data tape:
\begin{equation}
\mathcal{T} = \langle \mathcal{L}, c, \mathcal{R} \rangle \in \text{Peano}^* \times \text{Peano} \times \text{Peano}^*
\end{equation}
Pointer movements ($\langle, \rangle$) shift Peano numbers between the stacks $\mathcal{L}$ and $\mathcal{R}$ in $\mathcal{O}(1)$ type-level operations, and loops are evaluated via tail-recursive AST interpretation.

\section{Operational Semantics of Compilers}

\subsection{TypeScript: Structural Conditional Types}
In TypeScript, type functions evaluate via structural pattern matching:
\begin{lstlisting}[language=Java]
type Evolve<T, S extends number, Acc extends unknown[] = []> =
  Acc['length'] extends S
    ? T
    : Evolve<StepTape<T>, S, [...Acc, unknown]>;
\end{lstlisting}
TypeScript's type checker (\texttt{checker.ts}) performs memoized type instantiation. When the type in the recursive position is structurally naked and returns directly from the branch, the compiler activates tail-call recursion optimization.

\subsection{Rust: Nominal Horn-Clause Trait Resolution}
In Rust, types are nominal terms and trait bounds represent logic program clauses:
\begin{align}
\text{LocalRule}(L, C, R) &\to \text{StepCell}(L, C, R) \\
\text{Step}(Cons(N, Tail), P, C) &\leftarrow \text{Step}(Tail, C, N) \\
\text{Evolve}(T, Succ(K)) &\leftarrow \text{Evolve}(\text{Step}(T), K)
\end{align}
Rust resolves traits using Chalk-style SLD-resolution with coinductive cycle caches.

\subsection{C++20: Template Metaprogramming}
In C++20, computation is expressed via non-type template parameters and struct specialization:
\begin{lstlisting}[language=C++]
template<typename CurrentTape, int Steps>
struct Evolve {
  using type = typename Evolve<
    typename StepTape<CurrentTape>::type, Steps - 1>::type;
};
template<typename CurrentTape>
struct Evolve<CurrentTape, 0> {
  using type = CurrentTape;
};
\end{lstlisting}

\section{Circuit-Breaker Taxonomy}

Table~\ref{tab:circuit_breakers} summarizes the circuit breakers uncovered across the three languages.

\begin{table*}[t]
\centering
\caption{Taxonomy of Compiler Circuit Breakers across TypeScript, Rust, and C++}
\label{tab:circuit_breakers}
\begin{tabular}{@{}lllll@{}}
\toprule
\textbf{Compiler} & \textbf{Circuit Breaker Type} & \textbf{Default Limit} & \textbf{User Configurable?} & \textbf{Diagnostic Emitted} \\ \midrule
TypeScript 7.0.2 & Non-TCO Call-Stack Frame Fuse & 48 frames & No (Hardcoded) & \texttt{error TS2589} \\
TypeScript 7.0.2 & TCO Tail-Call Fuel Counter & 999 iterations & No (Hardcoded) & \texttt{error TS2589} \\
TypeScript 7.0.2 & Global Cumulative Instantiation Ceiling & $5,034,231$ instantiations & No (Hardcoded) & \texttt{error TS2589} \\
TypeScript 7.0.2 & Union Constituent Complexity Limit & $\sim 100,000$ terms & No (Hardcoded) & \texttt{error TS2590} \\
Rust 1.97.0 & Trait Goal Evaluation Recursion Depth & 128 frames (trips at 127) & Yes (\texttt{\#![recursion\_limit]}) & \texttt{error[E0275]} \\
Rust 1.97.0 & Chalk Coinductive Cycle Cache & Unbounded / Heuristic & No & Unhandled Hang / Freeze \\
Apple Clang C++20 & Template Instantiation Recursion Depth & 1024 frames (trips at 1024) & Yes (\texttt{-ftemplate-depth=N}) & \texttt{fatal error: max depth exceeded} \\ \bottomrule
\end{tabular}
\end{table*}

\section{Empirical Benchmarks and Findings}

\subsection{The TypeScript Tri-Fuse Hierarchy}
Our Phase 1 and Phase 2 experiments proved that TypeScript does not possess a single depth monitor, but a multi-tier hierarchy:
\begin{enumerate}
    \item \textbf{The Stack Frame Fuse ($D = 48$):} Without tail-call position, frame allocation is strictly capped at 48. At $D = 49$, compilation aborts immediately.
    \item \textbf{The Tail-Call Fuel Fuse ($F = 999$):} With accumulator tail-recursion, the compiler eliminates stack frames but increments a fuel counter. At $S = 999$, compilation succeeds ($551,145$ instantiations, $397.7$\,MB heap). At $S = 1000$, the compiler aborts with TS2589.
    \item \textbf{Quadratic Instantiation Scaling:} Instantiation growth scales quadratically:
    \begin{equation}
    \text{Inst}(S) = 0.4955 S^2 + 19.1858 S + 37327.42 \quad (R^2 = 0.999985)
    \end{equation}
\end{enumerate}

\subsection{Breadth Vulnerabilities and RAM Saturation}
By shifting from linear recursion to binary and quaternary branching trees ($b \in \{2, 4\}$), we demonstrated that depth limiters fail to protect compiler memory:
\begin{itemize}
    \item At depth $D = 16$ (well below the 48 stack limit), TypeScript allocated $4,359,612$ instantiations and saturated $2.16$\,GB of RAM in $6.5$\,s.
    \item At depth $D = 17$, the compiler hit the global ceiling of $5,034,645$ instantiations and aborted with TS2589.
    \item In quaternary trees, TypeScript heap surged to $2.84$\,GB at depth 9 in $13.56$\,s.
\end{itemize}

\subsection{Logarithmic Bypass: Trampolined Chunking}
We proved that TypeScript's tail-call fuel counter is dynamically scoped to the contiguous chain of a single type alias invocation. By introducing a trampolined chunking operator:
\begin{lstlisting}[language=Java]
export type StepChunk512<T> = EvolveTCO<T, 512>;
export type Trampoline512<T, Chunks, C = []> =
  C['length'] extends Chunks ? T :
  StepChunk512<T> extends infer NextT ?
    Trampoline512<NextT, Chunks, [...C, unknown]> : never;
\end{lstlisting}
Each chunk executes 512 steps within a safe envelope ($512 < 999$). The pattern \texttt{extends infer NextT} forces eager tuple resolution, resetting the fuel counter upon each bounce.
As a result, we scaled Rule 110 evolution from $S = 1,024$ ($2^{10}$) up to $S = 131,072$ ($2^{17}$) in \textbf{214\,milliseconds} with constant $\mathcal{O}(1)$ memory ($\approx 153.5$\,MB), shattering the 999-step ceiling by two orders of magnitude.

\subsection{Pathological Freeze: Subverting Cycle Detection}
Compilers rely on cycle detection heuristics to terminate cyclic searches. We constructed a 25-line pathological Rust program:
\begin{lstlisting}[language=C]
pub struct Nil;
pub struct Cons<H, T>(PhantomData<(H, T)>);
pub trait Eval { type Out; }
impl Eval for Nil { type Out = Nil; }
impl<H: Eval, T: Eval> Eval for Cons<H, T> {
  type Out = Cons<<H as Eval>::Out, <T as Eval>::Out>;
}
type T0 = Nil;
type T1 = Cons<T0, T0>; // ... T23 = Cons<T22, T22>
\end{lstlisting}
Because each sub-goal expands into syntactically unique type expressions, Rust's Chalk cycle cache never detects a repeating goal key. At depth 23 ($8,388,608$ nodes), \texttt{rustc} spent \textbf{30.39 seconds} in pure trait projection before successfully completing with exit code 0, completely evading circuit breakers.

\subsection{The C++20 Triad Performance Race}
Table~\ref{tab:triad_results} presents the comparative evaluation of Rule 110 evolution across Apple Clang C++20, Rust 1.97.0, and TypeScript 7.0.2.

\begin{table}[ht]
\centering
\caption{Triad Compiler Telemetry Benchmark ($W = 8$)}
\label{tab:triad_results}
\begin{tabular}{@{}rrrr@{}}
\toprule
\textbf{Steps ($S$)} & \textbf{TypeScript (ms)} & \textbf{Rust (ms)} & \textbf{Apple Clang (ms)} \\ \midrule
10 & 123.0 & 31.0 & \textbf{20.0} \\
100 & 151.0 & 35.2 & \textbf{28.7} \\
500 & 279.0 & 94.3 & \textbf{32.0} \\
1000 & 620.0 (TS2589) & 266.0 & \textbf{37.3} \\
1024 & FAIL & 275.0 & FAIL (Depth 1024) \\
10000 & 185.0 (Tramp) & 4820.0 & \textbf{481.9} \\ \bottomrule
\end{tabular}
\end{table}

Apple Clang emerges as the definitive speed champion:
\begin{itemize}
    \item At $S = 1000$, Clang resolves the cellular automaton in \textbf{37.3\,ms} (30\,ms real time) with $35.4$\,MB RSS—\textbf{8.8$\times$ faster than Rust} ($266$\,ms) and \textbf{20.6$\times$ faster than TypeScript} ($620$\,ms).
    \item When configured with \texttt{-ftemplate-depth=30000}, Clang computes 10,000 steps in \textbf{481.9\,ms} using $104.8$\,MB RSS.
\end{itemize}

\section{Architecture Recommendations}
Accidental Turing-completeness cannot be safely regulated by one-dimensional recursion limits. We propose three principles for future type system design:
\begin{enumerate}
    \item \textbf{Multi-Dimensional Thermodynamic Fuel Metering:} Compilers must measure total work volume $\mathcal{W} = \int (\text{depth} \times \text{breadth}) \, dt$ and heap consumption rates rather than simple call depths.
    \item \textbf{Hierarchical Lexical Scope Budgets:} Intermediate alias boundaries must not reset the global fuel budget; child evaluations must draw from an inherited quota.
    \item \textbf{Total Type-Level Sub-Dialects:} Non-domain metaprogramming should be confined to total functional sub-dialects (e.g., G\"odel's System $T$) where normalization is mathematically guaranteed.
\end{enumerate}

\section{Conclusion}
Project Chimera has established an exhaustive empirical and formal mapping of accidental Turing-completeness across TypeScript, Rust, and C++. We proved the existence of the TypeScript tri-fuse hierarchy, demonstrated the multi-gigabyte vulnerability of breadth branching, engineered a logarithmic trampoline bypass to 131,072 steps, induced a 30.4-second clean freeze in \texttt{rustc}, and established Clang C++20 as the compile-time execution champion. All source code, verification proofs, and empirical datasets are openly provided.

\bibliographystyle{IEEEtran}
\begin{thebibliography}{10}
\bibitem{turing1936}
A.~M. Turing, ``On computable numbers, with an application to the {Entscheidungsproblem},'' \emph{Proc. London Math. Soc.}, vol.~42, no.~2, pp. 230--265, 1936.

\bibitem{cook2004}
M.~Cook, ``Universality in elementary cellular automata,'' \emph{Complex Systems}, vol.~15, no.~1, pp. 1--40, 2004.

\bibitem{curry1958}
H.~B. Curry and R.~Feys, \emph{Combinatory Logic}, North-Holland, Amsterdam, 1958.

\bibitem{kleene1938}
S.~C. Kleene, ``On notation for ordinal numbers,'' \emph{The Journal of Symbolic Logic}, vol.~3, no.~4, pp. 150--155, 1938.

\bibitem{girard1989}
J.-Y. Girard, Y.~Lafont, and P.~Taylor, \emph{Proofs and Types}, Cambridge University Press, 1989.

\bibitem{chalk2020}
Rust Chalk Working Group, ``Chalk: An implementation of the Rust trait system based on Horn clauses,'' \url{https://github.com/rust-lang/chalk}, 2020.
\end{thebibliography}

\end{document}
